\section{Extensions to Two-Group Settings}\label{app:extensions}

The HDI+ROPE precision-goal framework presented in the main text for single-group
Bernoulli data generalises to two-group comparisons.
This appendix outlines the key adaptations for two-group binomial and continuous
outcomes, together with the corresponding sample-size planning equations.
For single-group continuous data, the expected stop iteration is derived in
Appendix~\ref{app:expected_stop}.
Empirical validation for utility of ePitG in these extensions is left for future work.

\subsection*{Two-Group Comparisons}

For two-group comparisons --- whether binomial ($\theta_1 - \theta_2$) or continuous
($\mu_1 - \mu_2$) --- the standard error of the difference serves as the precision
metric, enabling sequential A/B testing under the same HDI+ROPE decision framework
\citep{kruschke2013}.
% \citet{kruschke2013}: Kruschke's BEST paper introduces the HDI+ROPE framework for
% two-group comparisons of proportions and means, making it the primary reference for
% extending the HDI+ROPE decision rule to the two-group setting discussed here.
For the binomial case with conjugate Beta posteriors, the posterior of the difference
can be evaluated by direct i.i.d.\ sampling from each Beta distribution separately,
without Markov Chain Monte Carlo (MCMC) \citep{gelman2013bayesian}.
% \citet{gelman2013bayesian}: BDA3 (Chapter 2/3) is the standard reference for
% deriving posterior quantities by direct Monte Carlo simulation. Cited here because
% it supports drawing i.i.d. samples from independent Beta posteriors to obtain the
% difference distribution exactly, without requiring a Markov chain.
For the continuous case, under a weakly informative prior the posterior of
$\mu_1 - \mu_2$ converges to the frequentist sampling distribution by the
Bernstein--von Mises theorem \citep{gelman2013bayesian}, so the standard error from
Welch's approximation provides a reliable surrogate for the posterior HDI width when
samples are of moderate size or larger.
% \citet{gelman2013bayesian}: BDA3 also covers the Bernstein--von Mises theorem and
% the conditions under which posteriors converge to the frequentist sampling distribution
% under weakly informative priors (Chapter 4), justifying the use of Welch's standard
% error as a surrogate for the posterior HDI width.

\subsubsection*{Two-Group Binomial}

With independent Beta posteriors for $\theta_1$ and $\theta_2$, the HDI width of
the difference $\theta_1 - \theta_2$ is approximated via the standard error
\begin{equation}\label{eq:se_diff_binomial}
    \text{SE}(\hat\theta_1 - \hat\theta_2) \approx
    \sqrt{\frac{\theta_1(1-\theta_1)}{N_1} + \frac{\theta_2(1-\theta_2)}{N_2}}.
\end{equation}
This uses the CLT form with $N_i$ in the denominator; the exact Beta posterior
variance replaces $N_i$ with $N_i+1$ (cf.\ the Binomial Case in
Appendix~\ref{app:expected_stop}).
Setting the HDI width $2z_* \cdot \text{SE} = \omega_{\rm goal}$ and writing
$N_1 = rN_2$ for group-size ratio $r = N_1/N_2 > 0$, solving for $N_2$ gives
\begin{equation}\label{eq:ngoal_two_group_binomial}
    N_2 \approx \frac{4 z_{*}^{2}}{\omega_{\rm goal}^{2}}
        \left(\frac{\theta_1(1-\theta_1)}{r} + \theta_2(1-\theta_2)\right),
    \quad N_1 = r N_2.
\end{equation}
For equal group sizes ($r=1$, $N_1=N_2=N$) this simplifies to
\begin{equation}\label{eq:ngoal_two_group_binomial_equal}
    N \approx \frac{4 z_{*}^{2}[\theta_1(1-\theta_1)+\theta_2(1-\theta_2)]}{\omega_{\rm goal}^{2}}.
\end{equation}

\subsubsection*{Two-Group Continuous}

Following Welch's approximation, the standard error of the difference of means is
\begin{equation}\label{eq:se_diff_continuous}
    \text{SE}(\hat\mu_1 - \hat\mu_2) \approx
    \sqrt{\frac{\sigma_1^2}{N_1} + \frac{\sigma_2^2}{N_2}}.
\end{equation}
Writing $N_1 = rN_2$ for group-size ratio $r = N_1/N_2 > 0$ and solving for $N_2$ gives
\begin{equation}\label{eq:ngoal_two_group_continuous}
    N_2 \approx \frac{4 z_{*}^{2}}{\omega_{\rm goal}^{2}}
        \left(\frac{\sigma_1^2}{r} + \sigma_2^2\right),
    \quad N_1 = r N_2,
\end{equation}
with $\sigma_i^2$ replaced by $s_i^2$ when the population variances are unknown.
For equal group sizes ($r=1$, $N_1=N_2=N$) this simplifies to
\begin{equation}\label{eq:ngoal_two_group_continuous_equal}
    N \approx \frac{4 z_{*}^{2}(\sigma_1^2+\sigma_2^2)}{\omega_{\rm goal}^{2}}.
\end{equation}
